\section{Introduction}
\label{sec:introduction}

The companion paper~\citep[Revealed~Sacrifice]{lasser2026paper8a} proves the
\emph{Aggregate B-to-C Lower Bound} theorem for the GFM framework's exercised
capability volume $\volR$:
\[
  \volR^{U,[W]} \;\geq\; \sum_n \Delta\volL(X_n),
\]
where the sum runs over admissible trade events in a window $W$ whose bundles
are pairwise poset-independent.  The theorem converts the otherwise-unbounded
B-to-C gap (the divergence between $\volL$ capability possession and $\volR$
capability exercise) into a bounded one on the subset of capability-space
covered by observed sacrifice.

The theorem's admissibility conditions are load-bearing. Part~A requires S0
(valuation bounded by exercise), S1 (free choice with \emph{non-degrading
alternatives}), S2 (benchmark grounding), S4(a) (valuation additivity), and the
genuine-exchange conditions of the revealed-sacrifice event definition.  Part~B
additionally requires S3, S5, and within-bundle poset-independence.  Of these,
S1 is the most empirically delicate: not every observed sacrifice is voluntary.
An agent paying above market rate for medicine because they are in severe pain,
paying above market rate for water because they are dehydrated, or paying above
market rate for shelter because they are exposed are not voluntary trades in
the S1 sense.  Such sacrifices have \emph{reversed polarity}: more sacrifice
means the agent is in a \emph{worse} situation, not revealing stronger
preference.

This paper contributes the structural architecture that supplies a sufficient
filter condition for one below-sufficiency failure mode of S1
(Definition~\ref{def:s1_sufficiency}): events in which retention degrades the
agent and no third non-degrading alternative exists fail S1 and are filtered
from the aggregate.  Below-sufficiency \emph{targeting} where a third
non-degrading alternative is available is not by itself an S1 failure mode
(S1 may still hold via that third alternative), but classification requires
institutional attestation to confirm the third alternative's existence and
non-degrading status.  Other S1 failure modes proper (coercion, duress,
information asymmetry) likewise require institutional attestation rather
than this structural filter
(Remark~\ref{rem:s1_above}).  The paper also develops the diagnostic layers
that interpret what the lower bound is signalling when a low $\betaL$ fires an
alarm.

\subsection{What this paper does}

\begin{enumerate}
\item \textbf{Need-sufficiency architecture
  (Section~\ref{sec:need_sufficiency}).}
  A poset-construction precondition distinguishing needs from
  wants via multi-dimensional sufficiency thresholds.  The
  sufficiency threshold is the structural
  \emph{polarity boundary}: it identifies the region in
  which the H1+H2 below-sufficiency S1-failure mode operates
  (retention degrades the agent below threshold and no third
  non-degrading alternative exists), and the region in which
  the structural failure mode does not apply, leaving the
  remaining S1 conditions (coercion, duress, information
  asymmetry, etc.) to institutional attestation.  Events the
  attestation layer admits as S1-valid enter the companion
  paper's
  aggregate~\citep[Revealed~Sacrifice]{lasser2026paper8a};
  below-sufficiency targeting with a third non-degrading
  alternative likewise routes through attestation to determine
  S1-admissibility.  A
  polarity-correct benchmark inverts cost dimensions; a
  multiplicative downstream-gate on $\volL$-aggregation makes
  need-satisfaction a structural precondition for registering
  downstream weight; a parallel
  \emph{need-access cost} diagnostic tracks the benchmarked
  capability surrendered for need-satisfaction.

\item \textbf{Gap decomposition
  (Section~\ref{sec:gap_decomposition}).}
  A five-cell partition of the B-to-C gap---restricted,
  covered, dormant, residual, boundary-residual---with an exact
  identity $1 - \betaL = \sum \delta$ under the default
  literal-partition variant (strategy (c)), preserving
  seed-cell attribution so that $\delta_{\mathrm{covered}}$
  remains interpretable as proxy-failure on observed
  capabilities.  A four-cell variant under cooperative-closure
  semantics (strategies (a)/(b)) is also available, at the cost
  of priority-extension absorbing dormant and residual atoms
  into the covered cell.  Supplies a
  near-linear $\widetilde{O}(|P|+E+M)$ computability
  proposition from observed trade data plus
  governance-supplied structural inputs.  Interprets alarm firings: low $\betaL$ with high
  $\delta_{\mathrm{covered}}$ deficit is proxy failure on
  observed capabilities; low $\betaL$ with high
  $\delta_{\mathrm{restricted}}$ or $\delta_{\mathrm{residual}}$
  reflects structural invisibility rather than proxy failure.

\item \textbf{Wireheading-consistent concentration signal
  (Section~\ref{sec:wireheading}).}
  The Herfindahl--Hirschman Index on trade-flow leverage
  provides a tractable, governance-thresholded concentration
  alarm that is Schur-convex in the majorization order.  A
  Part-A category-flow variant is third-party observable from
  public committed events plus the public S1-admissibility
  labelling.

\item \textbf{Residual class characterization
  (Section~\ref{sec:residual_class}).}
  A revealed-observation carve of the residual class: the set
  of capabilities structurally outside the sacrifice channel's
  reach.  This differs from the classical
  benchmarkability-based residual class; an explicit witness
  shows $\ResB \setminus \ResS$ is non-empty, while
  $\ResS \setminus \ResB$ and $\ResB \cap \ResS$ depend on
  individuation choices
  (Proposition~\ref{prop:residual_intersection}).

\item \textbf{Worked example
  (Section~\ref{sec:worked_example}).}
  A two-substrate (biology + silicon) collective traced through
  four phases---healthy trade, automation-driven trade-flow
  concentration, diagnosis and recovery, and complementary-
  diagnostic cross-checking---demonstrates the full architecture
  end-to-end.
\end{enumerate}

\subsection{Relationship to the companion paper}
\label{sec:companion_relationship}

The companion paper~\citep[Revealed~Sacrifice]{lasser2026paper8a} establishes
the measurement theorem; this paper supplies the filter that keeps the channel
clean and the diagnostic architecture that interprets its output.  The division
of labor:

\begin{center}
\small
\begin{tabular}{@{}ll@{}}
\toprule
\textbf{Companion paper} (\citep[Revealed~Sacrifice]{lasser2026paper8a})
  & \textbf{This paper} \\
\midrule
Sacrifice observation model, S0--S5
  & Need-sufficiency architecture
  (S1 filter) \\
Aggregate B-to-C lower bound
  & Gap decomposition \\
Commitment-layer composition
  & Wireheading-consistent alarm \\
Two sacrifice channels
  & Residual class characterization \\
Non-self-balancing finding
  & Worked example \\
\bottomrule
\end{tabular}
\end{center}

We refer to the companion paper's definitions
(revealed-sacrifice event, assumptions S0--S5, the aggregate
B-to-C lower bound theorem, committed event, privacy theorem)
without re-deriving them; the companion paper's notation
conventions carry over.  \emph{Submission-package note.}  This
paper is intended for submission as a companion in the same
package as~\citep[Revealed~Sacrifice]{lasser2026paper8a}; reviewers should have
access to both.  Where a venue requires a single self-contained
submission, the load-bearing definitions (revealed-sacrifice
event, S0--S5, aggregate lower bound) would be inlined into a
unified single-paper version.  Stable archival identifiers
will accompany the sequence as it reaches archival status.

\subsection{Dependencies on prior papers}
\label{sec:dependencies}

This paper inherits the GFM
framework~\citep{lasser2026gfm,lasser2026poset,lasser2026horizon,lasser2026scm,lasser2026exo,lasser2026phase,lasser2026wamura}
and the measurement-theoretic infrastructure of the companion
paper~\citep[Revealed~Sacrifice]{lasser2026paper8a}.  A summary dependency table is
collected in Appendix~\ref{app:notation}.
